\documentclass[11pt]{article}
\usepackage{amsmath,amssymb}
\usepackage[hidelinks]{hyperref}
\title{An idle fault label does not certify restart}
\date{}
\begin{document}
\maketitle

A monitor can report no active fault because a second fault has
put the first component out of use. This is not evidence that
the first fault cleared. A restart also needs to respect work
left on the shared resource. One way to shorten the restart is
to validate the repaired component off the shared resource while
the old work settles. That extra fixture helps only when its
delay is smaller than the time it can save.

\section{What the published benchmark can establish}
The \href{https://drops.dagstuhl.de/entities/document/10.4230/OASIcs.DX.2026.7}
{HSU TwinFlow paper} evaluates anomaly detection and diagnosis using
a digital twin. It describes system reconfiguration as future work
and warns that the recorded full set of measurements need not be
available to a real controller. We checked the repository's actual
\href{https://github.com/imb-hsu/twinflow/blob/d8c623bbf3866b6edef4ef3f3b160229fd5c46f4/scripts/prepare_data.py}
{fault-label preparation}, Git blob
\texttt{6dde97dff447ee6f6d8fdb0e4c22423014a091b8}. The routine zeros
fault labels while the affected component is inactive, fills missing
fault fields with zero, and normalizes the nominal damping value.
These operations serve the benchmark's label semantics. Zero is not
a certificate that a previously faulty device has been repaired.

In an authored replay, feed the exact routine a mobile unit which
becomes idle for two readings and then becomes active. A retained
fault, an actual reset and missing raw values give the same
three-record prefix through the idle period: \((1,0,0)\).
When it is active again the retained fault gives \(1\) and the
other two examples give \(0\). This is a result about the source
transformation, not a claim about any particular released run.
It does not say that the complete measurements of a released
run would be insufficient under every sensor assignment.
The fixture shows exactly why neither a prepared zero nor a
reset timestamp may stand in for a fresh proof of health.
Do not feed hidden injection labels or future samples to the
runtime controller to make that test pass.

As a separate source check, the current
\href{https://j-ehrhardt.github.io/hai-cpps-benchmark/}
{HAI-CPPS documentation} explicitly separates online fields
from offline oracle states, useful practice for qualification.
The v2.2 campaign documented on its development branch enables
only one fault proxy per faulty recording. This is not an
independent compound recovery evaluation. The development
campaign description must not be attributed to an older
downloaded release. No test recordings or data-dependent
holdout selection were used here.

\section{A restart with an outstanding effect}
Consider an authored civil fixture in which a restarted motorised
cart must still finish a delivery. A drive issue was reported
before a subsequent obstruction disabled its active route.
An older accepted move physically occupies the shared test
lane. A reset is attempted while the cart is inactive. At the
decision boundary the same monitor prefix can arise with or
without a remaining drive issue. The two histories can both
also have the same reset acknowledgement. Restoring motion on
the lane is lawful only after actual old lane clearance and
a successful test of the drive. This is not a requirement to
put an unsafe probe on an occupied lane.

Use the same receipt types in a separate authored pump example:
the workload is to restore flow, the legacy effect is unsafe
pressure occupying the service envelope, and a no-load spin test
can instead be performed on a recirculation fixture before old
pressure relief is recognized. The examples stipulate the
relevant event semantics; neither asserts that such a fixture
is available on a particular plant. The pump test must in fact
validate the drive for the admitted service load; a no-load
test that is blind to permitted load-sensitive faults is not a
valid adapter. The two examples share the timed evidence rule,
but are not independently selected or externally validated
missions.

Keep four cuts distinct. \(t\) is when the component enters
the repair envelope in which a new test can succeed; an actual
test still must be witnessed. \(w\) is when the separate
inhibit from the second fault has cleared for the tester.
The local test instrument can start at \(x=\max(t,w)\) given
the corresponding enabling observations. \(y=s+c\) is when
actual settlement of old work at time \(s\) is established at
the relevant interlock, after a delay \(c\). Set \(c=0\) if
that interlock can witness clearance immediately. Let \(u>0\)
be the time of a successful shared-resource validation, \(h>0\)
the service duration, and \(d\ge0\) the extra elapsed time to
perform the same validation on a genuinely separate fixture
and return its receipt. Both modes use the same actual source
law, deadlines, observation scope and test validity. In the
second mode the added equipment, setup, test and receipt return
must be charged. Assume the two validations establish the same
predicate over the stated deployment envelope and that neither
is invalidated by another fault before issue. No precomputed
certificate decodes readiness from the suppressed fault label.

With no isolated fixture, validation starts only after both
\(x\) and \(y\), and the earliest finish is
\[
 S=\max(x,y)+u+h .
\]
An isolated validation starts at \(x\); it does not occupy the
old lane or vent and changes neither the ongoing duty nor its
release time. The new job starts only once both test and actual
settlement receipts are effective, so its earliest finish is
\[
 I=\max(x+u+d,y)+h .
\]
Each lower bound follows from the indicated causal predecessors.
Beginning each admissible step at its first available cut
attains the bound. These are exact within the stated test
and resource interface, not impossibility bounds for all
controllers with physical sensors or alternative repair
mechanisms. For a robust deadline apply the formula on each
actual admitted timing tuple rather than combine independent
worst cases if the times are constrained together.

The isolated fixture saves time exactly if \(d<y-x\).
When it saves time the savings are \(\min(u,y-x-d)\).
To see this, if \(y\le x\) no waiting can be overlapped.
Otherwise \(S=y+u+h\); subtract the larger of the
isolated test-ready time and release. In the replay's
example \(x=3,y=7,u=2,d=1,h=1\); only the isolated
fixture meets deadline 8, while shared validation needs
10. If the old work clears soon, extra test setup can
make isolation strictly worse. The direct same-resource
event scheduler attains exactly the same bounds. The
physical independence of the test is the intervention,
not superiority of the evidence rule over a direct scheduler.

\section{Acceptance boundary}
For a label-only rule the persistent and cleared source
prefixes are observationally identical. A true recovery can
admit service once valid evidence exists; a policy that treats
both identical idle observations as cleared can admit an
unsafe restart. Waiting for an effect-side validation is a
different observation contract and may be unable to meet
the original deadline on the same resources. Treat a
projection as sufficient only when the relevant histories
are separated at the actual decision cut, not merely after
reactivation or in simulator ground truth. An offline source
trace locates a fault, but is neither an authorization to
issue a replacement nor proof that the older job physically
settled. A live repair must satisfy both in the same run.

The replay exercises the exact upstream transformation on
synthetic inputs and compares the two scheduling formulas
against a discrete-event implementation on authored inputs.
It qualifies a source and gives a mechanism benchmark.
It does not test a holdout, close the integrated architecture
obligation, or change the enclosed-mission work.
\end{document}

